\begin{table}[!htbp]
\normalsize
\renewcommand{\arraystretch}{1.15}
\centering
\caption{Effect on migration intentions: pooled Spanish exposure, Logit average marginal effects}
\label{tab:post_pooled_share_ame_stata}
\setlength{\tabcolsep}{4pt}
\captionsetup{justification=centering}
\begin{tabular}{lcccc}
\hline
 & (1) & (2) & (3) & (4) \\
\hline
Spanish share 1936--1978 \(\times\) Post & 0.016* & 0.017** & 0.010 & 0.012 \\
 & (0.008) & (0.008) & (0.011) & (0.012) \\
\hline
Observations &        7,503 &        7,448 &        7,503 &        7,448 \\
Year FE & Yes & Yes & Yes & Yes \\
Municipality FE & Yes & Yes & Yes & Yes \\
Controls & No & Yes & No & Yes \\
Spanish share census & 1970 & 1970 & 1980 & 1980 \\
\hline
\end{tabular}
\par\vspace{1.5mm}
\begin{minipage}{0.98\linewidth}
\footnotesize\setlength{\parindent}{0pt}\setlength{\parskip}{1pt}
Notes: The dependent variable is migration intention. Each column reports a separate Logit specification. Entries are Logit average marginal effects on the 0--1 probability scale for \(Spanish\ share_{1936-1978} \times Post\). Spanish share is the share of Spanish-born immigrants who arrived between 1936 and 1978 over total municipal population, measured using the 1970 census and the 1980 census. All models include year and municipality fixed effects. Standard errors are clustered at the municipality level and reported in parentheses (56 clusters). Controls include age and male. * \(p<0.10\), ** \(p<0.05\), *** \(p<0.01\).
\end{minipage}
\end{table}
